% =====================================================================
% PROTOCOLO DE APLICACAO - EstufaNet
% Como usar no Overleaf:
%  - Para testar sozinho: compile este arquivo como esta.
%  - Para colar no relatorio: copie apenas o trecho entre
%    "INICIO DA SECAO" e "FIM DA SECAO", e adicione ao preambulo do
%    relatorio os pacotes e o \lstset abaixo.
% =====================================================================
\documentclass[12pt,a4paper]{article}
\usepackage[utf8]{inputenc}
\usepackage[T1]{fontenc}
\usepackage[brazilian]{babel}
\usepackage[margin=2.5cm]{geometry}
\usepackage{booktabs}
\usepackage{longtable}
\usepackage{array}
\usepackage{listings}
\usepackage{xcolor}
\usepackage[load-configuration=abbreviations]{siunitx}
\sisetup{output-decimal-marker={,}}

% --- Configuracao dos blocos de codigo (JSON/HTTP) ---
\lstset{
  basicstyle=\ttfamily\footnotesize,
  breaklines=true,
  frame=single,
  backgroundcolor=\color{gray!8},
  inputencoding=utf8,
  extendedchars=true,
  literate={á}{{\'a}}1 {é}{{\'e}}1 {í}{{\'i}}1 {ó}{{\'o}}1 {ú}{{\'u}}1
           {ã}{{\~a}}1 {õ}{{\~o}}1 {ç}{{\c c}}1 {â}{{\^a}}1 {ê}{{\^e}}1
           {Á}{{\'A}}1 {É}{{\'E}}1,
  columns=fullflexible,
  keepspaces=true
}

\begin{document}

% ================== INICIO DA SECAO ==================
\section{Protocolo de Aplica\c{c}\~ao}
\label{sec:protocolo}

\subsection{Vis\~ao geral}

O EstufaNet define um protocolo de aplica\c{c}\~ao pr\'oprio, sobre dois
mecanismos de transporte de aplica\c{c}\~ao que compartilham a mesma porta
TCP 8000 do servidor:

\begin{itemize}
  \item \textbf{HTTP/1.1}, para opera\c{c}\~oes de consulta, configura\c{c}\~ao e
        entrega da p\'agina do painel (modelo requisi\c{c}\~ao/resposta);
  \item \textbf{WebSocket} (RFC 6455), para o fluxo cont\'inuo de leituras em
        tempo real (conex\~ao persistente e bidirecional).
\end{itemize}

No WebSocket, todas as mensagens s\~ao \emph{frames} de texto contendo um
objeto JSON codificado em UTF-8. O protocolo \'e orientado a mensagens e
cada mensagem possui o campo obrigat\'orio \texttt{tipo}, que define sua
sem\^antica. Os pap\'eis s\~ao: \textbf{sensor} (cliente que produz leituras),
\textbf{painel} (cliente web que consome leituras) e \textbf{servidor}.

\subsection{Pontos de acesso (endpoints)}

\begin{table}[h]
\centering
\small
\begin{tabular}{@{}llp{7.2cm}@{}}
\toprule
\textbf{Tipo} & \textbf{Endere\c{c}o} & \textbf{Fun\c{c}\~ao} \\
\midrule
WebSocket & \texttt{/ws} & Canal dos sensores: envio de leituras e recebimento de \texttt{ack} ou \texttt{erro}. \\
WebSocket & \texttt{/ws/painel} & Canal do painel: recebe leituras repassadas pelo servidor. \\
HTTP GET  & \texttt{/} & Verifica\c{c}\~ao de status do servidor. \\
HTTP GET  & \texttt{/painel} & Entrega a p\'agina HTML do painel. \\
HTTP GET  & \texttt{/historico} & Consulta o hist\'orico de leituras gravadas. \\
HTTP GET  & \texttt{/alertas} & Consulta os alertas gravados. \\
HTTP GET  & \texttt{/limiar} & Consulta o limiar de temperatura vigente. \\
HTTP PUT  & \texttt{/limiar} & Altera o limiar de temperatura. \\
\bottomrule
\end{tabular}
\caption{Pontos de acesso do servidor (porta TCP 8000).}
\label{tab:endpoints}
\end{table}

\subsection{Mensagens WebSocket}

\subsubsection{\texttt{leitura} (sensor $\rightarrow$ servidor)}

Enviada pelo sensor a cada segundo, no canal \texttt{/ws}.

\begin{table}[h]
\centering
\small
\begin{tabular}{@{}llp{7.5cm}@{}}
\toprule
\textbf{Campo} & \textbf{Tipo} & \textbf{Regra} \\
\midrule
\texttt{tipo}      & texto  & Obrigat\'orio; deve ser exatamente \texttt{"leitura"}. \\
\texttt{device\_id}& texto  & Obrigat\'orio; identifica o dispositivo (ex.: \texttt{"sensor1"}). \\
\texttt{token}     & texto  & Obrigat\'orio; deve corresponder ao token cadastrado para o \texttt{device\_id}. \\
\texttt{temp}      & n\'umero & Obrigat\'orio; temperatura em \si{\celsius}, no intervalo $[-40,\,80]$. \\
\texttt{umid}      & n\'umero & Obrigat\'orio; umidade relativa em \%, no intervalo $[0,\,100]$. \\
\bottomrule
\end{tabular}
\caption{Campos da mensagem \texttt{leitura}. Valores booleanos n\~ao s\~ao aceitos como n\'umero. Campos extras s\~ao ignorados (o experimento de \emph{throughput} usa um campo \texttt{pad} apenas para aumentar o tamanho da mensagem).}
\end{table}

\begin{lstlisting}
{"tipo": "leitura", "device_id": "sensor1", "token": "tk-sensor1",
 "temp": 27.3, "umid": 61.2}
\end{lstlisting}

\subsubsection{\texttt{ack} (servidor $\rightarrow$ sensor)}

Resposta a toda leitura v\'alida. Confirma o recebimento e o armazenamento, e
informa se a leitura gerou alerta.

\begin{table}[h]
\centering
\small
\begin{tabular}{@{}llp{7.5cm}@{}}
\toprule
\textbf{Campo} & \textbf{Tipo} & \textbf{Descri\c{c}\~ao} \\
\midrule
\texttt{tipo}   & texto & Sempre \texttt{"ack"}. \\
\texttt{status} & texto & Sempre \texttt{"ok"}. \\
\texttt{alerta} & booleano & \texttt{true} se \texttt{temp} for estritamente maior que o limiar vigente. \\
\texttt{ts}     & texto & Instante de recebimento no servidor, formato ISO 8601 (ex.: \texttt{2026-11-03T12:36:22}). \\
\bottomrule
\end{tabular}
\caption{Campos da mensagem \texttt{ack}.}
\end{table}

\begin{lstlisting}
{"tipo": "ack", "status": "ok", "alerta": false, "ts": "2026-11-03T12:36:22"}
\end{lstlisting}

\subsubsection{\texttt{erro} (servidor $\rightarrow$ sensor)}

Enviada quando a mensagem recebida \'e rejeitada. A conex\~ao \textbf{permanece aberta}
ap\'os o erro, e o sensor pode enviar a pr\'oxima mensagem.

\begin{lstlisting}
{"tipo": "erro", "codigo": 401, "mensagem": "device_id ou token inválido"}
\end{lstlisting}

\subsubsection{\texttt{leitura} repassada (servidor $\rightarrow$ painel)}

Para cada leitura aceita, o servidor envia a todos os pain\'eis conectados em
\texttt{/ws/painel} um evento com o \texttt{tipo} \texttt{"leitura"} acrescido de
\texttt{ts} e \texttt{alerta}. O painel apenas recebe; qualquer texto enviado
por ele \'e lido e descartado (serve para manter a conex\~ao aberta).

\begin{lstlisting}
{"tipo": "leitura", "device_id": "sensor1", "temp": 33.8, "umid": 58.0,
 "ts": "2026-11-03T12:36:19", "alerta": true}
\end{lstlisting}

\subsection{C\'odigos de erro}

O servidor valida cada mensagem em ordem fixa: formato JSON, autentica\c{c}\~ao e
campos. Na primeira falha, responde com o c\'odigo correspondente e descarta a
mensagem (nada \'e gravado).

\begin{table}[h]
\centering
\small
\begin{tabular}{@{}cp{5.2cm}p{5.5cm}@{}}
\toprule
\textbf{C\'odigo} & \textbf{Causa} & \textbf{Exemplo} \\
\midrule
400 & Mensagem n\~ao \'e JSON v\'alido, ou n\~ao \'e um objeto JSON. & \texttt{abc}, \texttt{[1,2]} \\
401 & \texttt{device\_id} ausente/desconhecido ou token incorreto. & token \texttt{"errado"} \\
422 & \texttt{tipo} diferente de \texttt{"leitura"}, ou \texttt{temp}/\texttt{umid} ausente, n\~ao num\'erico ou fora da faixa. & \texttt{"temp": 200} \\
\bottomrule
\end{tabular}
\caption{C\'odigos de erro do canal WebSocket (reaproveitam a numera\c{c}\~ao do HTTP).}
\end{table}

\subsection{Sequ\^encia de opera\c{c}\~oes}

\paragraph{Sensor.}
\begin{enumerate}
  \item O sensor abre uma conex\~ao TCP com o servidor (porta 8000) e realiza o
        \emph{handshake} WebSocket: \texttt{GET /ws} com
        \texttt{Upgrade: websocket}, e o servidor responde
        \texttt{101 Switching Protocols}.
  \item O sensor envia uma mensagem \texttt{leitura}.
  \item O servidor valida (400, 401, 422). Se v\'alida, grava a leitura no banco,
        verifica o limiar, grava o alerta se houver, repassa o evento aos
        pain\'eis e responde \texttt{ack}. Se inv\'alida, responde \texttt{erro}.
  \item O sensor aguarda a resposta (o que permite medir o RTT da aplica\c{c}\~ao),
        espera 1~s e repete a partir do passo 2.
  \item Se a conex\~ao cair, o servidor registra a desconex\~ao em log; o sensor
        espera 3~s e tenta reconectar.
\end{enumerate}

\paragraph{Painel.}
\begin{enumerate}
  \item Ao abrir \texttt{GET /painel}, o navegador recebe o HTML e consulta por HTTP
        \texttt{GET /historico}, \texttt{GET /alertas} e \texttt{GET /limiar}.
  \item O painel abre o WebSocket \texttt{/ws/painel} e passa a receber as leituras
        em tempo real. Se a conex\~ao cair, tenta reconectar a cada 3~s.
  \item Para alterar o limite, o painel envia \texttt{PUT /limiar}.
\end{enumerate}

\begin{lstlisting}
Sensor                         Servidor                      Painel
  | --- GET /ws (Upgrade) ------> |                            |
  | <-- 101 Switching Protocols - |                            |
  | --- leitura ----------------> | valida, grava, alerta      |
  |                               | --- leitura (evento) ----> |
  | <-- ack (ou erro) ----------- |                            |
  |   ... a cada 1 s ...          |                            |
\end{lstlisting}

\subsection{Opera\c{c}\~oes HTTP}

\paragraph{\texttt{GET /historico}} Par\^ametros opcionais:
\texttt{device\_id} (filtra por dispositivo) e \texttt{limite} (inteiro de 1 a 1000,
padr\~ao 50). Devolve uma lista JSON, da leitura mais recente para a mais antiga.

\begin{lstlisting}
GET /historico?device_id=sensor1&limite=2 HTTP/1.1

HTTP/1.1 200 OK
[{"id": 812, "device_id": "sensor1", "temp": 27.3, "umid": 61.2,
  "ts": "2026-11-03T12:36:22"},
 {"id": 810, "device_id": "sensor1", "temp": 24.1, "umid": 55.0,
  "ts": "2026-11-03T12:36:20"}]
\end{lstlisting}

\paragraph{\texttt{GET /alertas}} Par\^ametro opcional \texttt{limite} (1 a 1000, padr\~ao 50).
Cada item cont\'em \texttt{id}, \texttt{device\_id}, \texttt{temp}, \texttt{limiar} (valor vigente no
momento do alerta) e \texttt{ts}.

\paragraph{\texttt{GET /limiar} e \texttt{PUT /limiar}} O corpo do \texttt{PUT} \'e um
objeto JSON com o campo num\'erico \texttt{valor}. Ambos respondem
\texttt{\{"temp": <limiar>\}}.

\begin{lstlisting}
PUT /limiar HTTP/1.1
Content-Type: application/json

{"valor": 25}

HTTP/1.1 200 OK
{"temp": 25.0}
\end{lstlisting}

Em requisi\c{c}\~oes HTTP com par\^ametros ou corpo inv\'alidos (por exemplo,
\texttt{limite=0} ou \texttt{valor} n\~ao num\'erico), o servidor responde
\texttt{422 Unprocessable Entity}, com detalhes no corpo.

\subsection{Regras de neg\'ocio}

\begin{itemize}
  \item \textbf{Alerta:} gerado quando \texttt{temp} $>$ limiar (compara\c{c}\~ao estrita). O limiar
        inicial \'e \SI{30.0}{\celsius}; \'e global (vale para todos os sensores) e fica em mem\'oria,
        voltando ao padr\~ao quando o servidor reinicia.
  \item \textbf{Autentica\c{c}\~ao:} cada dispositivo possui um token fixo
        (\texttt{sensor1}~a~\texttt{sensor5}, token igual a \texttt{tk-} mais o nome).
  \item \textbf{Persist\^encia:} banco SQLite com as tabelas \texttt{leituras}
        (\texttt{id}, \texttt{device\_id}, \texttt{temp}, \texttt{umid}, \texttt{ts}) e \texttt{alertas}
        (\texttt{id}, \texttt{device\_id}, \texttt{temp}, \texttt{limiar}, \texttt{ts}).
  \item \textbf{Registros (logs):} todo evento relevante (conex\~ao, leitura, alerta,
        erro, desconex\~ao, mudan\c{c}a de limiar) \'e escrito com \emph{timestamp}
        no console e no arquivo \texttt{estufanet.log}.
\end{itemize}

\subsection{Limita\c{c}\~oes conhecidas do protocolo}

\begin{itemize}
  \item O tr\'afego n\~ao \'e cifrado (\texttt{ws://} e \texttt{http://}, sem TLS); tokens e
        leituras trafegam em texto claro, o que \'e vis\'ivel nas capturas do Wireshark.
  \item Os tokens s\~ao fixos e est\~ao no c\'odigo-fonte; em produ\c{c}\~ao seriam
        gerados por dispositivo e armazenados de forma segura.
  \item As rotas HTTP n\~ao exigem autentica\c{c}\~ao: qualquer cliente com acesso \`a porta
        8000 pode ler o hist\'orico e alterar o limiar.
  \item O protocolo n\~ao possui numera\c{c}\~ao de mensagens nem reenvio na camada de
        aplica\c{c}\~ao; a confiabilidade \'e delegada ao TCP. Se o sensor cair antes do
        \texttt{ack}, n\~ao h\'a como saber se a leitura foi gravada.
\end{itemize}
% =================== FIM DA SECAO ===================

\end{document}
